\subsection{THE COMPLETION OF A UNIFORM SPACE}

THEOREM 3. Let X be a uniform space. Then there exists a complete Hausdorff uniform space \(\hat{\mathbf{X}}\) and a uniformly continuous mapping \(i: \mathbf{X} \to \hat{\mathbf{X}}\) having the following property:

\begin{enumerate}
\def\labelenumi{(\Alph{enumi})}
\setcounter{enumi}{15}
\tightlist
\item
  Given any uniformly continuous mapping \(f\) of \(\mathbf{X}\) into a complete Hausdorff uniform space \(\mathbf{Y}\) , there is a unique uniformly continuous mapping \(g: \hat{\mathbf{X}} \to \mathbf{Y}\) such that \(f = g \circ i\) .
\end{enumerate}

If \((i_1, X_1)\) is another pair consisting of a complete Hausdorff uniform space \(X_1\) and a uniformly continuous mapping \(i_1: X \to X_1\) having the property (P), then there is a unique isomorphism \(\varphi: \hat{X} \to X_1\) such that \(i_1 = \varphi \circ i\) .

The first statement of the theorem signifies that the pair \((i, \hat{\mathbf{X}})\) is the solution of the universal mapping problem (Set Theory, Chapter IV, § 3, no. 1) in which the \(\Sigma\) -sets are complete Hausdorff uniform spaces, the \(\sigma\) -morphisms are uniformly continuous mappings and the \(\alpha\) -mappings are uniformly continuous mappings of X into a complete Hausdorff uniform space. The uniqueness of the pair \((i, \hat{\mathbf{X}})\) up to a unique isomorphism therefore follows from the general properties of solutions of universal mapping problems (loc. cit.). It remains to prove the existence of the pair \((i, \hat{\mathbf{X}})\) .

\begin{enumerate}
\def\labelenumi{\arabic{enumi})}
\tightlist
\item
  Definition of \(\hat{\mathbf{X}}\) . Let \(\hat{\mathbf{X}}\) be the set of minimal Cauchy filters (no. 2) on \(\mathbf{X}\) . We shall define a uniform structure on \(\hat{\mathbf{X}}\) . For this purpose, if \(\mathbf{V}\) is any symmetric entourage of \(\mathbf{X}\) , let \(\tilde{\mathbf{V}}\) denote the set of all pairs \((\mathfrak{X}, \mathfrak{Y})\) of minimal Cauchy filters which have in common a V-small set. We shall show that the sets \(\tilde{\mathbf{V}}\) form a fundamental system of entourages of a uniform structure on \(\hat{\mathbf{X}}\) :
\end{enumerate}

\begin{enumerate}
\def\labelenumi{(\roman{enumi})}
\item
  Since each \(\mathcal{X} \in \hat{\mathbf{X}}\) is a Cauchy filter, we have by definition \((\mathcal{X}, \mathcal{X}) \in \tilde{\mathbf{V}}\) for every symmetric entourage \(\mathbf{V}\) of \(\mathbf{X}\) ; hence axiom \((\mathbf{U}_{\mathbf{I}}^{\prime})\) is satisfied.
\item
  If V and \(V'\) are two symmetric entourages of X, then \(W = V \cap V'\) is a symmetric entourage, and every set which is W-small is also V-small and \(V'\) -small; hence \(\tilde{W} \subset \tilde{V} \cap \tilde{V}'\) , which prove \((B_{I})\) .
\item
  The sets \(\tilde{\mathbf{V}}\) are symmetric by definition, hence \((\mathbf{U}_{\mathbf{II}}^{\prime})\) is satisfied. (iv) Given a symmetric entourage \(\mathbf{V}\) of \(\mathbf{X}\) , let \(\mathbf{W}\) be a symmetric entourage such that \(\tilde{\mathbf{V}} \subset \mathbf{V}\) . Consider three minimal Cauchy filters \(\mathcal{X}, \mathcal{Y}, \mathfrak{Z}\) such that \((\mathcal{X}, \mathcal{Y}) \in \tilde{\mathbf{W}}\) and \((\mathcal{Y}, \mathfrak{Z}) \in \tilde{\mathbf{W}}\) ; then there are two W-small sets \(\mathbf{M}, \mathbf{N}\) such that \(\mathbf{M} \in \mathcal{X} \cap \mathcal{Y}\) and \(\mathbf{N} \in \mathcal{Y} \cap \mathfrak{Z}\) . Since \(\mathbf{M}\) and \(\mathbf{N}\) belong to \(\mathcal{Y}\) , \(\mathbf{M} \cap \mathbf{N}\) is not empty and therefore (no. 1, Proposition 1) \(\mathbf{M} \cup \mathbf{N}\) is \(\tilde{\mathbf{W}}\) -small and hence V-small; since \(\mathbf{M} \cup \mathbf{N}\) belongs to \(\mathcal{X}\) and to \(\mathfrak{Z}\) we have \(\tilde{\mathbf{W}} \subset \tilde{\mathbf{V}}\) ; hence \((\mathbf{U}_{\mathbf{III}}^{\prime})\) is satisfied.
\end{enumerate}

We show next that the uniform space \(\hat{\mathbf{X}}\) is Hausdorff. Let \(\mathcal{X},\mathcal{Y}\) be two minimal Cauchy filters on \(\mathbf{X}\) such that \((\mathcal{X},\mathcal{Y})\in \hat{\mathbf{X}}\) for all symmetric entourages \(\mathbf{V}\) of \(\mathbf{X}\) . It follows immediately that the sets \(\mathbf{M}\cup \mathbf{N}\) , where \(\mathbf{M}\in \mathcal{X}\) and \(\mathbf{N}\in \mathcal{Y}\) , form a base of a filter \(\mathfrak{Z}\) coarser than \(\mathcal{X}\) and \(\mathcal{Y}\) . Now \(\mathfrak{Z}\) is a Cauchy filter, since for every symmetric entourage \(\mathbf{V}\) of \(\mathbf{X}\) there is by hypothesis a V-small set \(\mathbf{P}\) belonging to both \(\mathcal{X}\) and \(\mathcal{Y}\) and therefore belonging to \(\mathfrak{Z}\) . By the definition of minimal Cauchy filters, we have \(\mathcal{X} = \mathfrak{Z} = \mathcal{Y}\) , and this shows that \(\hat{\mathbf{X}}\) is Hausdorff.

\begin{enumerate}
\def\labelenumi{\arabic{enumi})}
\setcounter{enumi}{1}
\item
  Definition of \(i\) ; the uniform structure of \(\mathbf{X}\) is the inverse image under \(i\) of that of \(\hat{\mathbf{X}}\) . We know that for each \(x \in \mathbf{X}\) the neighbourhood filter \(\mathfrak{B}(x)\) of \(x\) in \(\mathbf{X}\) is a minimal Cauchy filter (no. 2, Proposition 5, Corollary 1). So we define \(i(x) = \mathfrak{B}(x)\) . Let \(f = i \times i\) ; we shall show that for each symmetric entourage \(\mathbf{V}\) of \(\mathbf{X}\) we have \(\overline{j}(\tilde{\mathbf{V}}) \subset \mathbf{V} \cup \overline{j}[(\mathbf{V})^{\sim}]\) , and this will prove our assertion (§ 2, no. 4). Now, if \([i(x), i(y)] \in \tilde{\mathbf{V}}\) , there is a V-small set \(\mathbf{M}\) which is a neighbourhood of each of \(x\) and \(y\) , hence \((x, y) \in \mathbf{V}\) . Conversely, if \((x, y) \in \mathbf{V}\) , it is immediately seen that the set \(\mathbf{V}(x) \cup \mathbf{V}(y)\) is \(\overline{\mathbf{V}}\) -small and is a neighbourhood of each of \(x\) and \(y\) .
\item
  \(\hat{\mathbf{X}}\) is complete and \(i(\mathbf{X})\) is dense in \(\hat{\mathbf{X}}\) . The trace on \(i(\mathbf{X})\) of a neighbourhood \(\tilde{\mathbf{V}}(\mathfrak{X})\) of a point \(\mathfrak{X} \in \hat{\mathbf{X}}\) is the set of all \(i(x)\) such that
\end{enumerate}

\[
(\mathcal {X}, i (x)) \in \tilde {\mathbf {V}}.
\]

This relation means that there is a V-small neighbourhood of \(x\) in X which belongs to \(\mathcal{X}\) , i.e.~that \(x\) is an interior point of a V-small set of \(\mathcal{X}\) . Let M be the union of the interiors of all V-small sets of \(\mathcal{X}\) ; then M belongs to \(\mathcal{X}\) (no. 2, Proposition 5, Corollary 4) and from what has been said it follows that \(\tilde{\mathrm{V}}(\mathcal{X}) \cap i(\mathrm{X}) = i(\mathrm{M})\) . We conclude that:

\begin{enumerate}
\def\labelenumi{(\roman{enumi})}
\item
  \(\tilde{\mathbf{V}}(\mathcal{X}) \cap i(\mathbf{X})\) is not empty, hence \(i(\mathbf{X})\) is dense in \(\hat{\mathbf{X}}\) .
\item
  The trace of \(\tilde{\mathbf{V}}(\mathcal{X})\) on \(i(\mathbf{X})\) belongs to the filter base \(i(\mathcal{X})\) on X; hence this filter base converges in \(\hat{X}\) to the point \(\mathfrak{X}\).
\end{enumerate}

Now let \(\mathfrak{F}\) be a Cauchy filter on \(i(\mathbf{X})\) ; then from 2) above and Proposition 4 of no. 1, \(\overline{i}^{1}(\mathfrak{F})\) is a base of a Cauchy filter \(\mathfrak{G}\) on \(\mathbf{X}\) . Let \(\mathfrak{X}\) be a minimal Cauchy filter coarser than \(\mathfrak{G}\) (no. 2, Proposition 5); then \(i(\mathfrak{X})\) is a Cauchy filter base on \(i(\mathbf{X})\) (no. 1, Proposition 3), and \(\mathfrak{F}=i[\overline{i}^{1}(\mathfrak{F})]\) is finer than the filter whose base is \(i(\mathfrak{X})\) . Since the latter converges in \(\hat{\mathbf{X}}\) , so does \(\mathfrak{F}\) , and Proposition 9 of no. 4 therefore shows that \(\hat{\mathbf{X}}\) is complete.

\begin{enumerate}
\def\labelenumi{\arabic{enumi})}
\setcounter{enumi}{3}
\tightlist
\item
  Verification of the property (P). Let \(f\) be a uniformly continuous mapping of \(\mathbf{X}\) into a complete Hausdorff uniform space \(\mathbf{Y}\) . Let us first show that there is a unique uniformly continuous mapping \(g_{0}: i(\mathbf{X}) \to \mathbf{Y}\) such that \(f = g_{0} \circ i\) . Since \(f\) is continuous, we have
\end{enumerate}

\[
f (x) = \lim f (\mathfrak {B} (x)),
\]

hence if we put \(g_0(i(x)) = \lim f(\mathfrak{V}(x))\) , we have \(f = g_0 \circ i\) ; so it remains to show that \(g_0\) is uniformly continuous in \(i(\mathbf{X})\) . Let \(\mathbf{U}\) be an entourage of \(\mathbf{Y}\) and let \(\mathbf{V}\) be a symmetric entourage of \(\mathbf{X}\) such that the relation \((x, x') \in \mathbf{V}\) implies \((f(x), f(x')) \in \mathbf{U}\) ; we have seen in 2) that the relation \((i(x), i(x')) \in \tilde{\mathbf{V}}\) implies \((x, x') \in \mathbf{V}\) , hence also is implied

\[
(g _ {0} (i (x)), g _ {0} (i (x ^ {\prime}))) \in \mathbf {U},
\]

which proves our assertion.

Let \(g\) be the extension of \(g_0\) by continuity to \(\hat{\mathbf{X}}\) (no. 6, Theorem 2); then \(f = g \circ i\) , and it is clear that \(g\) is the unique continuous mapping of \(\hat{\mathbf{X}}\) into \(Y\) satisfying this relation, since \(i(\mathbf{X})\) is dense in \(\hat{\mathbf{X}}\) (Chapter I, § 8, no. 1, Proposition 2, Corollary 1).

Q.E.D.

DEFINITION 4. The complete Hausdorff uniform space \(\hat{\mathbf{X}}\) defined in the proof of Theorem 3 is called the Hausdorff completion of \(\mathbf{X}\) , and the mapping \(i: \mathbf{X} \to \hat{\mathbf{X}}\) is called the canonical mapping of \(\mathbf{X}\) into its Hausdorff completion.

We note also the following facts:

PROPOSITION 12. (i) The subspace \(i(\mathbf{X})\) is dense in \(\hat{\mathbf{X}}\) .

\begin{enumerate}
\def\labelenumi{(\roman{enumi})}
\setcounter{enumi}{1}
\item
  The graph of the equivalence relation \(i(x) = i(x')\) is the intersection of the entourages of \(\mathbf{X}\) .
\item
  The uniform structure of X is the inverse image under i of that of \(\hat{X}\) {[}or of that of the subspace \(i(\mathbf{X})\) {]}.
\item
  The entourages of \(i(\mathbf{X})\) are the images under \(i \times i\) of the entourages of X, and the closures in \(\hat{X} \times \hat{X}\) of the entourages of \(i(\mathbf{X})\) form a fundamental system of entourages of \(\hat{X}\) .
\item
  and (iii) have been proved in the course of the proof of Theorem 3; (iv) is a consequence of (i) and (iii) by virtue of general results proved earlier (§ 2, no. 4, Remark and Proposition 6). The relation
\end{enumerate}

\[
i (x) = i \left(x ^ {\prime}\right)
\]

means by definition that \(x\) and \(x'\) have the same neighbourhood filter. But this implies, by definition, that \((x, x') \in V\) for every entourage \(V\) of \(X\) , and the converse is obvious.

COROLLARY. If X is a Hausdorff uniform space, then the canonical mapping \(i: X \to \hat{X}\) is an isomorphism of X onto a dense subspace of \(\hat{X}\) .

When X is Hausdorff, \(\hat{X}\) is said to be the completion of X, and we generally identify X with a dense subset of \(\hat{X}\) by means of i.

Remark. If this identification is made, the minimal Cauchy filters on X are just the traces on X of the neighbourhood filters of points of \(\hat{X}\) ; this follows from the proof of Theorem 3.

The Corollary to Proposition 12 characterizes the completion of a Hausdorff uniform space:

PROPOSITION 13. If Y is a complete Hausdorff uniform space and X a dense subspace of Y, then the canonical injection \(\mathbf{X} \to \mathbf{Y}\) extends to an isomorphism of \(\hat{\mathbf{X}}\) onto Y.

For every uniformly continuous mapping of X into a complete Hausdorff uniform space Z extends uniquely to a uniformly continuous mapping of Y into Z by Theorem 2 of no. 6.

PROPOSITION 14. Let X be a complete Hausdorff uniform space, U its uniformity, and let Z be a dense subspace of X. If U' is a uniformity on X which is coarser than U and which induces the same uniformity as U on Z, then U = U'.

Let \(X'\) denote the set \(X\) with the uniformity \(\mathcal{U}'\) . The composition of the canonical mapping \(X' \to \hat{X}'\) and the identity mapping \(X \to X'\) is a uniformly continuous mapping \(\varphi: X \to \hat{X}'\) . Since \(Z\) is Hausdorff for the uniform structure induced by \(\mathcal{U}'\) , the restriction of \(\varphi\) to \(Z\) is by hypothesis an isomorphism of \(Z\) onto the dense subspace \(\varphi(Z)\) of \(\hat{X}'\) ; it follows (no. 6, Corollary to Theorem 2) that \(\varphi\) itself is an isomorphism of \(X\) onto \(\hat{X}'\) , hence \(X' = \hat{X}'\) and \(\mathcal{U}' = \mathcal{U}\) .

PROPOSITION 15. Let X and X' be two uniform spaces. For each uniformly continuous mapping \(f: \mathbf{X} \to \mathbf{X}'\) there is a unique uniformly continuous mapping \(\hat{f}: \hat{\mathbf{X}} \to \hat{\mathbf{X}}'\) such that the diagram

\[
\begin{array}{c}\mathbf {X} \stackrel {{i ^ {\prime}}} {{\to}} \mathbf {X} ^ {\prime}\\i \downarrow \quad \quad \quad \quad \quad \quad \quad \quad \quad \quad \quad \quad \quad \quad \quad \quad \quad \quad \quad \quad \hat {\mathbf {X}} \stackrel {{\rightarrow}} {{\longrightarrow}} \hat {\mathbf {X}} ^ {\prime}\\f\end{array}
\]

is commutative (*), where \(i: \mathbf{X} \to \hat{\mathbf{X}}\) and \(i': \mathbf{X}' \to \hat{\mathbf{X}}'\) are the canonical mappings.

Apply Theorem 3 to the function \(i' \circ f: \mathbf{X} \to \hat{\mathbf{X}}'\) .

COROLLARY. If \(f: \mathbf{X} \to \mathbf{X}'\) and \(g: \mathbf{X}' \to \mathbf{X}''\) are two uniformly continuous mappings and \(h = g \circ f\) , then \(\hat{h} = \hat{g} \circ \hat{f}\) .

This is an immediate consequence of the uniqueness in Proposition 15.

